\section{Discussion}
\label{sec:discussion}

In the CS1 holdout (139 students), the discrimination AUC of the 30-second change, in the matched comparison, fell
within the band $[0.45, 0.55]$ set in advance (Section~\ref{sec:limitations} qualifies the error rate). Equivalence here
means that each signal, used alone as a ranking score, separated stall run-ups from matched control run-ups no better
than the band allows, on average over the pairs weighted within writers. It does not mean that typing does not change
before a stall, that stalls and controls did not both change, that a combination of signals could not separate them, or
that the signals are uninformative for every writer.

For research and practice that builds triggers for help from keystrokes, the result says that the two signals of this
estimator gave no basis for a trigger in the seconds before a stall. What a difference inside the band would mean for a
trigger can be read from the conversion on seen writers, made after the band was set (Appendix~\ref{app:threshold}): at
an AUC of 0.55, 8.9--10.0\% of the firings of a threshold trigger would fall on pauses that become stalls (8.2--8.9\% at
0.50), and even at 0.70 only 12.9--18.4\%. Those who plan to use an estimator of the same type as a trigger can check
the run-up signal against a band set in advance on held-out writers, which is the recommendation of the literature on
evaluation design \cite{gardner2018evaluating,lakens2017equivalence,lijmer1999empirical}; our data do not show the harm
of omitting the check.

\subsection{What was established and what remains possible}

What the data establish is limited to the statistic above being inside the band in the matched comparison, including
after removing stalls that contain events inside the pause. Three further readings remain possible and were not tested.
(a) Unmatched evaluation might make run-up rises look large. In these data the unmatched comparison showed no rise, and
in the learner description the matched comparison was higher than the unmatched windows, so the data neither support
nor contradict this reading, nor the opposite reading that matching removed a real run-up difference
\cite{mansournia2018case}. (b) Standard quantities and flexible learners might hold run-up differences that the
estimator's signals miss. The learner analysis is an independent description; in its main analysis the relation to the
upper end of the band was undecidable, and its sensitivity analyses moved in both directions. (c) The instrument might
have been weakened on this input. Its levels discriminated less in CS1 (0.524 and 0.521) than in the essays (0.569 and
0.580), the input was rebuilt from edit events with a synthesised release, and its coefficients were not estimated on
any of these data. The equivalence cannot be told apart from (c) (Section~\ref{sec:limitations}).

\subsection{The direction below 0.5}

All four point estimates lay slightly below 0.5. In the terms of equivalence testing
\cite{lakens2018equivalence}, three combinations were statistically different from 0.5 in the downward direction and
equivalent, and the Stuck-state probability with $\varepsilon = 0$ was equivalent and not different. The registration
did not ask about the direction. Candidate causes, none of them tested, are the exclusion that only controls receive
(controls close to a later onset are removed), the different shares at the floor (61\% and 65\%), and stalls with events
inside the pause (removing them gave 0.505 and 0.500); each moves the estimate within the band. In the essays the
direction was the reverse.

\subsection{Relation to theories of pausing}

Under the view that planning proceeds in parallel with typing, planning that cannot be completed in parallel appears in
the pause itself rather than in the typing before it \cite{roeser2025typing,olive2002concurrent}. These are claims about
language boundaries in prose, and applying them to long stops in code is an extrapolation. Our result is consistent
with this view but is not a test of it, because a serial view that allows planning to start late in the run-up would
predict the same.

\subsection{Calibration}

The calibration was needed in this study: the decision rule, applied to boundary configurations drawn from seen writers,
declared equivalence more often than its nominal level in one of its two versions (Section~\ref{sec:design}). It is
known for regression that the coverage of cluster-level inference can depart from its nominal level when clusters are
few or of unequal size \cite{cameron2008bootstrap,cameron2015practitioner,mackinnon2023cluster}. Our values are
consistent with that, by analogy, but we did not test the mechanism; designs with few effective writers may benefit from
a similar calibration.

\subsection{Reporting matched controls}

In other fields, unmatched controls have made accuracy look larger \cite{lijmer1999empirical,pepe2008pivotal}. When
reporting a run-up signal, stating what the comparison partner was matched on lets readers judge the result, although
our data do not show an effect of the way of matching.

\subsection{Materials}

We release the preregistration, the deviation, the start records and the results registration (all public on the Open
Science Framework), the code of the evaluation procedure (the parts that do not depend on the estimator), the list of
holdout writers as identifiers of the public dataset, the seed, the computational environment and the manifest of the
verified files with their SHA-256 values, the version and SHA-256 of the replay tool, and the steps to obtain the public
data. The public data themselves are not redistributed. Whether the estimator, the replay tool and the conversion
scripts, and the minimal estimate rows needed for the decision (the two signals at $u$ and $u - 30$~s and the pair
identifiers), can be released under the terms of the original data is to be settled before submission.
